\scanpage{25}
(où $f$ est identifié à un élément de $\unsure{\mathscr P|\mathfrak P}_3(S)$).

Cette bijection est évidemment compatible avec l'opération de $G_f\simeq D_3$ \inserted{$=\unsure{\mathfrak S|\mathscr S}_3$}, opérant sur \deleted{\unsure{$\mathbf Z/2$}} $\mathrm{Or}(f)$ via le quotient $\mathfrak S_2$ de $\mathfrak S_3$ (i.e. via signature) et sur $f$ \unsure{tautologique|tautologiqmt|tautologiquement}.

\deleted{Donc} [On retrouve que $G_f$ \deleted{conserve} \inserted{opère} par permutations paires sur $T$ --] On voit ainsi que \deleted{l'image de} $G_f$ \deleted{\unsure{dans $G$|dans}} \inserted{opère \unsure{fidèlt|fidèlement} sur $T$, et son image dans} $\mathfrak S_T$ (\deleted{\unsure{et est|contenue dans}} $\mathrm{Alt}_T$) est donc un ss-groupe d'ordre $6$. Donc l'image de $G$ dans $\mathrm{Alt}_T$ est d'ordre un multiple de $6$, et par les résultats \unsure{précédts|précédents} un multiple de $10$, donc un multiple de $30$.

\uline{Proposition} Soit $a\unsure{\in|=}\{s,t\}\in A$. Considérons \deleted{l'ens $\varphi$} les deux faces incidentes à $a$, et \inserted{l'ens $\varphi\in\mathscr P_2(S)$} leurs troisième\unsure{s} sommets $u,v$, \deleted{\unsure{Alors|dans $A$}} \deleted{de leurs arêtes} Considérons $a\times\varphi\longrightarrow A$ défini par $(s,u)\longmapsto\{s,u\}$. Alors les cinq éléments $\mathrm{tr}\,\overline{\{s,u\}}$ $(s\in a,u\unsure{\in|+}\varphi)$, $\mathrm{tr}\,\bar a$ de $T$ sont distincts.

Résulte des deux prop. précédentes.
% archive sheet number: 24
